\begin{lemma}
\label{lemma-derivation-extends}
Let $R$ be a ring. Let $D : R \to R$ be a derivation.
\begin{enumerate}
\item For any ideal $I \subset R$ the derivation $D$ extends
canonically to a derivation $D^\wedge : R^\wedge \to R^\wedge$
on the $I$-adic completion.
\item For any multiplicative subset $S \subset R$ the derivation
$D$ extends uniquely to the localization $S^{-1}R$ of $R$.
\end{enumerate}
If $R \subset R'$ is a finite type extension of rings such that
$R_g \cong R'_g$ for some $g \in R$ which is a nonzerodivisor in $R'$,
then $g^ND$ extends to $R'$ for some $N \geq 0$.
\end{lemma}

\begin{proof}
Proof of (1). For $n \geq 2$ we have $D(I^n) \subset I^{n - 1}$
by the Leibniz rule. Hence $D$ induces maps $D_n : R/I^n \to R/I^{n - 1}$.
Taking the limit we obtain $D^\wedge$. We omit the verification that
$D^\wedge$ is a derivation.

\medskip\noindent
Proof of (2). To extend $D$ to $S^{-1}R$ just set
$D(r/s) = D(r)/s - rD(s)/s^2$ and check the axioms.

\medskip\noindent
Proof of the final statement. Let $x_1, \ldots, x_n \in R'$ be generators
of $R'$ over $R$. Choose an $N$ such that $g^Nx_i \in R$.
Consider $g^{N + 1}D$. By (2) this extends to $R_g$. Moreover, by
the Leibniz rule and our construction of the extension above we have
$$
g^{N + 1}D(x_i) = g^{N + 1}D(g^{-N} g^Nx_i) = -Ng^Nx_iD(g) +
gD(g^Nx_i)
$$
and both terms are in $R$. This implies that
$$
g^{N + 1}D(x_1^{e_1} \ldots x_n^{e_n}) =
\sum e_i x_1^{e_1} \ldots x_i^{e_i - 1} \ldots x_n^{e_n} g^{N + 1}D(x_i)
$$
is an element of $R'$. Hence every element of $R'$ (which can be written
as a sum of monomials in the $x_i$ with coefficients in $R$) is mapped to an
element of $R'$ by $g^{N + 1}D$ and we win.
\end{proof}
